\chapter{Decoding with a state machine}
\MARSChapterHeading{Decoding with a state machine}
\label{revision:qoi}

PPM gave each pixel three explicit sample bytes. An image with a large uniform
background repeats those bytes even when the next pixel is predictable. \href{https://qoiformat.org/qoi-specification.pdf}{QOI},
the Quite OK Image format, represents pixels with short operations: repeat a
color, recover one seen earlier, describe a small change, or supply new components.
Decoding now requires remembering what previous operations established.

The complete \texttt{qoi-codec.asm} in \texttt{revision/labs/files} supplies a
bounded validator, decoder, encoder and file pipeline. Its worked image has
eight pixels and uses all six QOI chunk types. We will trace that image before
considering the larger sprite. The main problem is keeping the source cursor,
remaining pixel count, previous color and cache consistent at every transition.
The independent task changes output placement while preserving that transition
order.

This chapter assumes bounded parsing, RGB words, unsigned byte operations,
function calls and file ownership. Our local version-1.0 domain limits width
and height to one through 64. It accepts channels three or four and colorspace
zero or one, reconstructs all four components, and requires one exact image
with its complete end marker and no trailing bytes.

\section{Separate file structure from reconstructed pixels}

A QOI header has fourteen bytes: magic \texttt{qoif}, two big-endian unsigned
32-bit dimensions, a channels byte and a colorspace byte. The worked header
specifies width eight, height one, channels four and colorspace zero. The chunks
start at offset fourteen. An eight-byte end marker follows them: seven zero
bytes and a final one.

The eight pixels need 32 raw RGBA bytes. Their worked chunks need sixteen,
but header and marker add 22, making the file 38 bytes. A tiny file need not be
smaller than its raw samples; this fixture is chosen to exercise every chunk type.

Channels and colorspace are descriptive fields. Channels three does not prohibit
an RGBA chunk or authorize discarding its alpha. Colorspace does not ask this
codec to transform the component values. Preserve the flags in the metadata
and the reconstructed components in the output. Colors are not premultiplied:
a red value of seven and alpha 128 remain those separate values.

The local descriptor has four unsigned words: width, height, channels and
colorspace, in that order. These words use native little-endian storage, unlike
the dimensions in the serialized header. A framebuffer pixel is word
\texttt{0xAARRGGBB}; byte view therefore shows B, G, R, A.

\begin{figure}[htbp]
  \centering
  \begin{tabular}{l l}
    \toprule
    Representation           & The same reconstructed pixel \\
    \midrule
    RGBA components          & \texttt{07 06 05 80}         \\
    Framebuffer word         & \texttt{80070605}            \\
    Little-endian word bytes & \texttt{05 06 07 80}         \\
    RGB-only display         & Red 7, green 6, blue 5       \\
    \bottomrule
  \end{tabular}
  \caption{Alpha 128 survives decoding even when an RGB tool does not display it. The components are not premultiplied by alpha. A successful RGB preview alone cannot establish the stored alpha or component-byte order. Values in code font are hexadecimal.}
  \label{fig:revision-qoi-representations}
\end{figure}


A variable-size operation is called a chunk here. It is not a PNG chunk with a
length, type and CRC. A QOI opcode determines how many more bytes belong to its
operation. Keep that distinction when reusing the preceding chapter's parser
reasoning: the availability proof remains useful, but the record layout changes.

\section{Define the state before dispatching an opcode}

Begin with previous color $(R,G,B,A)=(0,0,0,255)$. Initialize every one of the
64 cache entries to $(0,0,0,0)$. These are different states: the initial predictor
is opaque black, while an untouched cache slot is transparent black. An initial
INDEX of slot zero therefore has a defined transparent result.

Track the next unread chunk byte and the number of pixels still required by
$W\cdot H$. Also retain the preceding INDEX slot, or a sentinel meaning that the
last chunk was not INDEX. Our validator rejects consecutive INDEX chunks naming
the same slot; repetition should use RUN. A non-INDEX chunk clears that history.
Cache contents alone cannot establish whether this sequence rule was followed.

\begin{table}[htbp]
  \centering
  \begin{tabular}{l l r l}
    \toprule
    Operation & Leading byte                             & Total bytes & Effect                            \\
    \midrule
    RGB       & \texttt{fe}                              & 4           & Set RGB; retain alpha             \\
    RGBA      & \texttt{ff}                              & 5           & Set all components                \\
    INDEX     & \texttt{00xxxxxx}                        & 1           & Recover cached color              \\
    DIFF      & \texttt{01xxxxxx}                        & 1           & Three small RGB changes           \\
    LUMA      & \texttt{10xxxxxx}                        & 2           & Green change plus residuals       \\
    RUN       & \texttt{11xxxxxx}, except \texttt{fe/ff} & 1           & Repeat previous color 1--62 times \\
    \bottomrule
  \end{tabular}
  \caption{Test the complete \texttt{fe} and \texttt{ff} tags before classifying the upper two bits. A chunk's byte count differs from its produced pixel count; RUN is the case where one input byte produces several pixels. Bit patterns use binary digits.}
  \label{tab:revision-qoi-tags}
\end{table}


Full-byte tags \texttt{fe} and \texttt{ff} take precedence. Testing the upper
two bits first would mistake them for runs of 63 and 64 pixels. The worked
walker dispatches them before extracting those bits:
\Needspace{11\baselineskip}
\MARSCodeLabel{Recognize full-byte tags first}
\begin{minted}{mars-idea}
li $t1, 254
beqc $t0, $t1, walk_rgb
nop
li $t1, 255
beqc $t0, $t1, walk_rgba
nop
dsrl $t1, $t0, 6
\end{minted}

The available chunk interval ends eight bytes before the selected file end.
Every opcode and operand read must stay before that boundary. Reserving the
marker first prevents its mapped zero bytes from completing a truncated RGB,
RGBA or LUMA operation. A pixel count and an input cursor have separate jobs:
one bounds reconstruction, the other bounds reads.

\section{Recover changes without losing a component}

RGB replaces red, green and blue while retaining the previous alpha. RGBA replaces
all four. INDEX copies all four components from its selected cache slot. DIFF
and LUMA retain alpha and reconstruct RGB with modulo-256 arithmetic.

DIFF packs three two-bit fields into its low six bits. Subtract two from each
field to recover a difference from $-2$ through one. Add that difference to the
corresponding previous component, then mask to eight bits. Field value three
means a change of one, not three. A change of $-2$ applied to one gives 255,
not a negative stored component and not a saturated zero.

The green part of the worked DIFF handler is:
\Needspace{9\baselineskip}
\MARSCodeLabel{Remove the bias and wrap the reconstructed component}
\begin{minted}{mars-idea}
dsrl $t8, $t0, 2
andi $t8, $t8, 3
daddiu $t8, $t8, -2
daddu $t5, $t5, $t8
andi $t5, $t5, 255
\end{minted}

Here \texttt{\$t0} is the opcode and \texttt{\$t5} initially contains the
previous green value. Mask after addition; masking a negative difference before
adding may be equivalent modulo 256, but it does not make that difference an
ordinary positive signed quantity. State which arithmetic interpretation your
range tests and component operations use.

LUMA spends six bits on the green difference and four bits each on the red and
blue residuals relative to it. From the first byte recover
$\Delta G=(\mathit{tag}\mathbin{\mathrm{and}}63)-32$. From the next byte recover residuals
$e_R=(\mathit{next}\gg4)-8$ and $e_B=(\mathit{next}\mathbin{\mathrm{and}}15)-8$.
Then add $\Delta G+e_R$ to red, $\Delta G$ to green and $\Delta G+e_B$ to blue,
wrapping each result to a byte.

For \texttt{a5 97}, $\Delta G=5$, $e_R=1$ and $e_B=-1$.
Following previous RGB $(1,1,1)$, this produces $(7,6,5)$.
The residual is not the full red or blue difference. Treating it as one loses
part of the shared change. This encoding exploits related channel movement;
it does not convert the pixel into a different stored color space.

\section{Update the cache and count every produced pixel}

After reconstructing a color, put it in slot
\[
  (3R+5G+7B+11A)\mathbin{\mathrm{and}}63.
\]
The weighted sum's low six bits select one of 64 entries. This is a replacement
cache, not a list of every earlier color. For example, opaque black and opaque
$(64,0,0)$ both hash to slot 53. Inserting the latter replaces the former.
An encoder must compare the complete cached color before selecting INDEX;
a matching slot number alone does not prove a matching color.

RUN repeats the previous color between one and 62 times. Its low six bits store
one less than the repetition count. Check that count against the remaining
pixels before emitting any of them. Runs can cross row boundaries; rows have
no special delimiter in the chunk stream.

RUN also inserts its color into the cache. This matters at the start of the
worked image: \texttt{c0} emits opaque black and installs it in slot 53. The
later byte \texttt{35}, an INDEX of decimal slot 53, recovers that color.
A decoder that skips the insertion would instead recover transparent black.
An RGB-only preview might make that error invisible.

\begin{figure}[htbp]
\centering
\begin{tabular}{r l l l r}
\toprule
File offset & Chunk bytes & Operation & Reconstructed RGBA & Cache slot\\
\midrule
14 & \texttt{c0} & RUN, one pixel & \texttt{00 00 00 ff} & 53\\
15 & \texttt{fe ff 00 00} & RGB & \texttt{ff 00 00 ff} & 50\\
19 & \texttt{c0} & RUN, one pixel & \texttt{ff 00 00 ff} & 50\\
20 & \texttt{35} & INDEX 53 & \texttt{00 00 00 ff} & 53\\
21 & \texttt{7f} & DIFF & \texttt{01 01 01 ff} & 4\\
22 & \texttt{a5 97} & LUMA & \texttt{07 06 05 ff} & 11\\
24 & \texttt{ff 07 06 05 80} & RGBA & \texttt{07 06 05 80} & 22\\
29 & \texttt{04} & INDEX 4 & \texttt{01 01 01 ff} & 4\\
\bottomrule
\end{tabular}
\caption{The worked eight-pixel stream uses every chunk type. Offsets 0--13 hold its header, and 30--37 hold its end marker. Cache slots are decimal; chunk and component bytes are hexadecimal. The first RUN inserts opaque black into slot 53, making the later INDEX meaningful.}
\label{fig:revision-qoi-trace}
\end{figure}


The remaining trace illustrates the other state dependencies. DIFF \texttt{7f}
changes black to $(1,1,1,255)$ and places it in slot four. LUMA changes RGB but
keeps alpha 255. RGBA then changes alpha to 128 without changing those RGB
values. The last INDEX restores all four components from slot four, including
alpha 255. It does not merely reuse the current alpha with older RGB.

Each emitted pixel reduces the remaining count by one. A RUN may do that
several times without consuming another opcode. When the count reaches zero,
require the cursor to equal the reserved marker boundary and verify all eight
marker bytes. Too few pixels, a run past the required extent, extra chunk bytes,
an altered marker and trailing data are distinct reasons to reject the stream.

\section{Validate completely before publishing output}

\texttt{qoi\_validate} receives full byte pointer, unsigned selected count and
an eight-byte-aligned descriptor pointer in \texttt{\$a0/\$a1/\$a2}. Reject a
count above 20502 before pointer access; counts below 23 are format errors.
For accepted dimensions, require exactly $W\cdot H$ reconstructed pixels and
the exact end marker before publishing the sixteen descriptor bytes. Success
returns status zero and pixel count in \texttt{\$v1}; a rejected stream leaves
the descriptor unchanged.

\texttt{qoi\_decode} receives source pointer, byte count, word destination and
unsigned pixel capacity. It validates the whole stream first, then checks
capacity, then publishes reconstructed words. Success returns packed dimensions,
height in the high 32 bits and width in the low 32 bits. Format failure wins over
capacity failure, and neither may write output. A final-byte error must not
leave an otherwise convincing prefix in the caller's framebuffer.

The implementation validates, then walks again to emit. Each walk starts with
a fresh predictor, zeroed cache, cursor and INDEX history; validation's final
state would give emission the wrong start. No state from an earlier call,
successful or failed, may seed the next image.

Two passes revisit valid bytes, but give this implementation its whole-output
failure guarantee without a private full-image staging buffer. A different
implementation may stage results privately or validate another way. Preserve
the published behavior when comparing designs; removing validation to obtain
a lower instruction count changes the interface being measured.

\section{Encode with legal alternatives and a proved bound}

The encoder receives a framebuffer pointer, descriptor pointer, byte destination
and unsigned capacity. It validates dimensions and flags before capacity, and
requires capacity at least $22+5WH$ before reading pixels or writing output.
Twenty-two accounts for header and marker; five bytes per pixel allows an RGBA
literal for every pixel. For the eight-pixel example the required capacity is
62 even though its worked encoded representation occupies only 38 bytes.

This is a conservative reservation contract, not a prediction of the final size.
It avoids partially emitting a header before discovering that a later pixel
needs a larger chunk. Success reports the exact emitted count. Padding, an
extra NUL and the educational write service's guard byte are outside that count.
At the maximum local dimensions the reservation is 20502 bytes.

An encoder may prefer a repeated-pixel run, an exact cache hit, DIFF, LUMA,
RGB or RGBA. Several legal streams can describe the same image. Choosing a
literal where another encoder chooses a difference is not itself an error.
Checks reconstruct dimensions, flags and all RGBA components rather than
matching the reference's chunk bytes.

The local interface nevertheless requires useful compression of a uniform
image. Opaque black needs at most $22+\lceil N/62\rceil$ bytes because it matches
the initial predictor. Any other uniform color needs at most
$27+\lceil(N-1)/62\rceil$: allow one five-byte literal, then runs. A shorter
cache-based first chunk is also valid. Flush a pending run both when it reaches
62 and when the image ends; a changed next color must not erase an unflushed run.

RGB may be used only when retaining the previous alpha is correct. Otherwise
use RGBA unless an exact INDEX or repeated-pixel operation already represents
all four components. Channels three is not permission to remove an alpha change.
The whole-operation checks include uniform bounds, cache collisions, changed
flags and repeated calls as well as ordinary pictures.

\section{Join decoding and encoding through real files}

The supplied caller offers bound, validate, decode, encode, load and save modes;
mode six joins load, decode, encode and save. The loader returns a byte count,
the decoder returns packed dimensions, and metadata retains descriptive flags
for the encoder. Each stage stops the chain on error. A returned pixel count
cannot substitute for a serialized byte count.

The file rules follow Chapter~\ref{revision:files}: establish EOF using actual
read counts, probe privately when the offered buffer is exactly full, close
before validating loaded bytes, and validate output before a truncating open.
Saving loops on positive short writes and fails on zero progress or an error.
Load may leave an accepted input prefix after failure; metadata remains unchanged.
A failed save after opening may leave a partial file, but must close its descriptor.

All six completed entries return status zero on success. Status one denotes a
format or descriptor error, two an unsupported count or capacity, and three file
I/O failure; every failure returns zero in \texttt{\$v1}. Bound returns one for
unsupported dimensions. Updates preserve full saved registers, frame pointer,
stack pointer and return path, with no console I/O. Byte/path pointers may be
unaligned; pixel pointers are four-byte aligned and descriptors eight-byte aligned.
Caller objects are disjoint.

A round trip through your own encoder and decoder can preserve the same mistake
in both. The lab therefore compares actual output files with an independent
host decoder and with the intended RGBA, not only with the original serialized
bytes. It also checks late errors and untouched output separately from successful
round trips. An attractive sprite preview is one observation among these checks.

\section{Mirror the image without reversing the state machine}

Complete \texttt{qoi\_decode\_mirrored} in \texttt{qoi-mirrored.asm}. Its arguments,
capacity units, status precedence and failure guarantees match \texttt{qoi\_decode}.
Return the same packed dimensions, but place reconstructed pixel $(x,y)$ at
$(W-1-x,y)$. Retain all four components exactly. There is no vertical reflection.

The compressed stream still has to be interpreted in its original order.
An INDEX depends on earlier cache insertions; a difference depends on the previous
pixel, and a run may cross rows. Reading opcode bytes backward does not produce
a reflected image and may not even preserve chunk boundaries.

One design decodes successfully with the supplied routine, then swaps words
within each output row. Another keeps the forward decoder state and changes
its emission address. For emitted pixel number $k$, use $y=\lfloor k/W\rfloor$
and $x=k\bmod W$, then address
\[
  destination+4\bigl(yW+(W-1-x)\bigr).
\]
A run needs a new placement for each produced pixel, including across rows.
A single reversed walk over the entire framebuffer would also reverse row order.

You may reuse the completed helpers. Write only the accepted pixel interval
on success and no output on failure; preserve source and other caller data.
The lab checks one-column, one-row, odd-width and multirow cases, full alpha,
state reuse and rejected inputs. Explain both the validation boundary and the
placement rule before drawing conclusions from the displayed image.
